\section{Introduction}\label{sec:intro}

Misleading information can affect financial decisions through attention,
beliefs and confidence. Evidence on manipulated financial news and social
media during bank runs establishes the relevance of information transmission
for financial markets \citep{kogan2023social,clarke2021fake,cookson2023social}.
Social-media amplification is not itself evidence that the amplified claims
were false. A distinct macroeconomic question is whether an information
environment characterized by frequent disinformation is associated with worse
growth outcomes, particularly when financial conditions are fragile.

This paper investigates that question with a country panel and an
illustrative financing mechanism. Its empirical contribution is a transparent
assessment of the information content of a broad disinformation proxy,
including conditional mean associations, quantile slopes, fragility
interactions and chronological forecast evaluation. It does not estimate
the causal macroeconomic effect of financial misinformation or the returns
to a particular content intervention.

The mechanism combines signal contamination with an explicitly assumed
noise-sensitive lending rule. The resulting investment shortfall has
increasing differences in contamination and fragility: deterioration in
information can matter more when financing constraints are close to binding.
A separate feedback example permits multiple steady states. A scalar
persistence calculation gives a conditional comparison of equal proportional
changes in model parameters. These results clarify hypotheses and identify
assumptions needed for policy analysis; they are not substitutes for
measurement or identification.

The empirical proxy combines five DSP indicators of government and party
disinformation in 175 countries over 2000--2025. It measures relative intensity
in the political information environment, rather than the share of financial
messages that are false. The analysis retains that distinction
throughout. The local projections estimate associations with cumulative
outcomes, while separate quantile regressions assess whether slopes differ
across the conditional growth distribution. Financial stress is a
conditioning state, not an externally assigned treatment.

The three-year cumulative-growth coefficient is
$\GrowthTwo$ percentage points per standard deviation of the index at mean
conditioning states. The corresponding banking-stress interaction is
$\FragilityTwo$ points. Two-year conditional growth slopes are
$\QuantileTen$, $\QuantileFifty$ and $\QuantileNinety$ points at the tenth,
fiftieth and ninetieth percentiles. These numbers must be evaluated with their
uncertainty, sample variation and reverse-timing diagnostics, rather than
interpreted as treatment effects. Joint bootstrap comparisons do not distinguish the lower-tail slope
from the median or upper tail, so the data do not support a
disproportionately large downside association.

The forecasting extension compares the same macro-financial benchmark with
and without the index. It fits component scaling and ranks solely within
each training window and respects a one-year predictor-availability lag.
Across held-out target years 2015--2024, the change in tenth-quantile pinball
loss is reported as a $\ForecastTail$ percent gain, where negative values
mean worse performance. The evaluated design does not support improved downside-risk
prediction from adding the index. Latest-vintage source data and a
short evaluation period remain material limitations.

\paragraph{Related literature.}
Models of misinformation and narrative formation study platform incentives,
social learning and competing explanations
\citep{acemoglu2024model,muellerfrank2024weakest,eliaz2020model}.
Experimental work documents motivated evaluation and incomplete updating after
retractions \citep{thaler2024fake,goncalves2026retractions}. Evidence on
political misperceptions and institutional confidence reinforces the need to
distinguish a political-information proxy from financial-content exposure
\citep{cariolle2024misinformation,guriev20213g,zhuravskaya2020political}.
Assessments of misinformation also warn against extrapolating from vivid
incidents to aggregate harm \citep{budak2024misunderstanding,altay2023misinformation}.

The macroeconomic motivation comes from narrative economics
\citep{shiller2017narrative,shiller2019narrative}, vulnerable-growth methods
\citep{adrian2019vulnerable}, and collateral-based amplification
\citep{kiyotaki1997credit,bernanke1999financial}. The feedback illustration
is related to uncertainty traps \citep{fajgelbaum2017uncertainty}.
The empirical tools are local projections \citep{jorda2005estimation} and
two-step panel quantile estimation \citep{canay2011simple}. Applying those
tools to a persistent observational index does not by itself identify a
structural shock. The paper's contribution lies in the combination of
measurement discipline, reproducible conditional evidence, and explicit
predictive validation.

Section~\ref{sec:model} develops the illustrative mechanism;
Section~\ref{sec:data} explains measurement; Section~\ref{sec:empirics}
presents estimation and validation. Sections~\ref{sec:risk} and
\ref{sec:policy} delimit risk and policy interpretations, and
Section~\ref{sec:conclusion} concludes. The appendix provides proofs,
robustness results and replication details.
